\section{Introduction}
\label{sec:intro}

Large language models answer medical questions fluently and score well on medical exam benchmarks \citep{singhal2023clinical}, yet they can state facts that no source supports, a failure known as hallucination \citep{ji2023hallucination}. In medicine a wrong dose or an outdated threshold reads exactly like a correct one. A model's knowledge also stops at its training date, while guidelines are revised, and it cannot say where a fact comes from.

Retrieval-augmented generation (RAG) addresses these problems by giving the model passages from a chosen document collection at the time of the question \citep{lewis2020rag}. The answer can then be traced to a passage, and the collection can be updated without retraining the model. RAG does not remove the risk by itself: retrieval always returns passages, even for a question the collection cannot answer, and the model may still answer from memory. A medical system must also show which passage supports each claim, and refuse when the collection does not hold the answer.

The project brief sets this target: answer medical questions from WHO guidelines and PubMed articles, cite the sources, and refuse when the information is missing. A refusal is acceptable; an unsupported answer is not.

\paragraph{Objectives.} We set five objectives:
\begin{enumerate}
\item Build the pipeline of the brief: 50 documents, 512-token chunks with 20\% overlap, BGE-M3 embeddings in Chroma, top-5 retrieval, and a Gradio app.
\item Make every answer cite the retrieved passages by number, and check the citations in code before the answer is shown.
\item Refuse with one fixed sentence when the passages do not answer the question, and measure how often the system refuses correctly and how often it refuses wrongly.
\item Compare the system with the two baselines the course requires: the model with no retrieval, and BM25 keyword retrieval.
\item Report the results with an error analysis, and state what the system cannot do.
\end{enumerate}

\paragraph{Contributions.} Beyond the brief's minimum: a BM25 written from scratch and fused with dense retrieval; a cross-encoder refusal gate chosen with no language model call; a citation check that turns an answer without a valid passage number into a refusal; a 40-question set with a held-out near-domain group; and an evaluation that fits a quota of 50 free requests a day.

\paragraph{Outline.} Sections~\ref{sec:related} to \ref{sec:setup} cover related work, the data, the method and the setup; Sections~\ref{sec:results} to \ref{sec:conclusion} the results, an error analysis, the discussion and the conclusion. The last section declares the tools used, including AI assistance.
